\subsection{Localization of flatness}

PROPOSITION 13. Let S be a multiplicative subset of a ring A and M an A-module. If M is flat (resp. faithfully flat), \(S^{-1}M\) is a flat (resp. faithfully flat) \(S^{-1}A\) -module and a flat A-module.

As \(S^{-1}M = M \otimes_{A} S^{-1}A\) , the first assertion follows from Chapter I, §2, no. 7, Corollary 2 to Proposition 8 (resp. Chapter I, §3, no. 3, Proposition 5); moreover, \(S^{-1}A\) is a flat A-module (§2, no. 4, Theorem 1); hence if M is a flat A-module, so is \(S^{-1}M\) by virtue of Chapter I, §2, no. 7, Corollary 3 to Proposition 8.

Remark. If N is an \(S^{-1}A\) -module, \(S^{-1}N\) is identified with N and this is consequently equivalent to saying that N is a flat \(S^{-1}A\) -module or a flat A-module.

PROPOSITION 14. Let A be a ring, B a commutative A-algebra and T a multiplicative subset of B. If N is a B-module which is flat as an A-module, \(T^{-1}N\) is a flat A-module.

\(T^{-1}N = T^{-1}B \otimes_{B} N\) ; the proposition then follows from Chapter I, § 2, no. 7, Proposition 8, applied with A replaced by B, B by A, E by \(T^{-1}B\) and F by N.

PROPOSITION 15. Let A, B be two rings, \(\phi: A \to B\) a homomorphism and N a B-module. The following properties are equivalent:

\begin{enumerate}
\def\labelenumi{(\alph{enumi})}
\item
  N is a flat A-module.
\item
  For every maximal ideal \(\mathfrak{n}\) of \(\mathbf{B}\) , \(\mathbf{N}_{\mathfrak{n}}\) is a flat A-module.
\item
  For every maximal ideal \(\mathfrak{n}\) of \(\mathbf{B}\) , if we write \(\mathfrak{m} = \bar{\phi}^{-1}(\mathfrak{n})\) , \(\mathbf{N}_{\mathfrak{n}}\) is a flat \(\mathbf{A}_{\mathfrak{m}}\) -module.
\end{enumerate}

For all \(a \notin m\) , the homothety of \(N_{n}\) induced by a is bijective, hence \(N_{n}\) is canonically identified with \((\mathrm{N}_{\mathrm{n}})_{\mathrm{m}}\) and the equivalence of (b) and (c) follows from the

Remark following Proposition 13; the fact that (a) implies (b) is a special case of Proposition 14. It remains to prove that (b) implies (a), that is, that, if (b) holds, for every injective A-module homomorphism \(u: \mathbf{M} \to \mathbf{M}'\) , the homomorphism \(v = 1 \otimes u: \mathbf{N} \otimes_{\mathbf{A}} \mathbf{M} \to \mathbf{N} \otimes_{\mathbf{A}} \mathbf{M}'\) is injective. Now, \(v\) is also a B-module homomorphism and, for it to be injective, it is necessary and sufficient that \(v_n: (\mathbf{N} \otimes_{\mathbf{A}} \mathbf{M})_n \to (\mathbf{N} \otimes_{\mathbf{A}} \mathbf{M}')_n\) be so for every maximal ideal \(n\) of \(B\) (no. 3, Theorem 1). As

\[
(\mathbf {N} \otimes_ {\mathbf {A}} \mathbf {M}) _ {n} = \mathbf {B} _ {n} \otimes_ {\mathbf {B}} (\mathbf {N} \otimes_ {\mathbf {A}} \mathbf {M}) = \mathbf {N} _ {n} \otimes_ {\mathbf {A}} \mathbf {M},
\]

\(v_{n}\) is just the homomorphism \(1 \otimes u: N_{n} \otimes_{A} M \to N_{n} \otimes_{A} M'\) , which is injective since \(N_{n}\) is a flat A-module by hypothesis.

COROLLARY. For an A-module M to be flat (resp. faithfully flat), it is necessary and sufficient that, for every maximal ideal m of A, \(M_{m}\) be a flat (resp. faithfully flat) \(A_{m}\) -module.

The necessity of the conditions follows from Proposition 13. Conversely, if \(M_{m}\) is a flat \(A_{m}\) -module for every maximal ideal m of A, M is a flat A-module by virtue of Proposition 15 applied to the case where \(\phi\) is the identity. Finally, if \(M_{m}\) is a faithfully flat \(A_{m}\) -module for all m, then \(mM_{m} = mA_{m}M_{m} \neq M_{m}\) , hence \(mM \neq M\) for all m (no. 3, Proposition 9), which proves that M is a faithfully flat A-module (Chapter I, § 3, no. 1, Proposition 1 (d)).

